\section{Threats to validity}
\label{sec:limitations}

This section is in two parts. \S\ref{sec:lim-named} names the eight limitations
that bound what \S\ref{sec:results} may be read to show, each with the experiment
that would resolve it and none of which has been run. \S\ref{sec:lim-validity}
groups the remaining threats by the kind of validity each bears on. Caveats
stated in full at the point of measurement are cross-referenced rather than
restated.

\subsection{Named limitations, and the experiment each needs}
\label{sec:lim-named}

\paragraph{L1. The strength gain is small, and often within noise.}
Head to head \vfive{} scores \numPanelVfiveVfour\% against \vfour{}
[\numStyleCIVfour] --- $\numDeltaVfour$ points on \vfour{}'s
\numPanelVfourVfour\% mirror baseline, on an interval that contains even money.
The nine-style means are level at \numVfiveMean\% against
\numVfourMean\%, and on minimax regret over that style set \vfour{} is
marginally the better of the two, \numVfourRegret{} against \numVfiveRegret{}.
Most of the per-style differences are a point or
less --- hunter $\numDeltaHunter$, bluffer $\numDeltaBluffer$, random
$\numDeltaRandom$, detective $\numDeltaDetective$, diversifier
$\numDeltaDiversifier$, \vthree{} $\numDeltaVthree$ --- and the largest,
\vtwo{} $\numDeltaVtwo$ and lockout $\numDeltaLockout$, are of a size the
per-cell intervals do not separate. No style in the table is one on which the
two policies' intervals fail to overlap. Nothing in this report should be read as
\vfive{} playing better than \vfour{}; the case for it is the pathology table
and the deception panel. \emph{Resolving experiment:} a paired common-deal
evaluation at an order of magnitude more deals per cell, with a per-style
optimum computed for each style, so that regret is measured against a best
response and not against the better of two arms \cite{xu-mirror}.

\paragraph{L2. No exploitability probe was run against \vfive{}.} The v0.4 study
ran one; this one did not, and the battery of \S\ref{sec:results} contains no
best-response search of any kind. \vfive{} therefore inherits no robustness
property from \vfour{}'s probe and claims none of its own. This is a strict
scope limit rather than a hedge: the deception panel measures performance against
three \emph{fixed} deceptive rules, which is not the same quantity as
performance against a policy that is searching for this policy's worst case.
\emph{Resolving experiment:} a local best-response probe against \vfive{},
reported as an achieved score that is a lower bound within its search class
\cite{lisy-lbr}, run with and without M6, since a deterministic argmax sits at
the deterministic leakage bound whatever its feature weights
\cite{alvim-leakage}.

\paragraph{L3. Mechanisms M3--M7 are specified and unbuilt, and two patches for
them are unmeasured.} \S\ref{sec:mech-designed} specifies the net-information ask
term, the per-seat knowledge model, target-dimension selection, partner-aware
stochastic selection and the online opponent model. None is on \vfive{}'s
decision path. Patches for M4/M5 and for M7 exist in the repository
under \filepath{research/v05/patches/}; neither has been evaluated in play, and
no number anywhere in this report is attributable to either. They are recorded
because unmeasured code that looks finished is a hazard, not because they support
a claim. \emph{Resolving experiment:} apply each patch, measure it alone against
the frozen configuration and then jointly, in the build order of
\S\ref{sec:mechanisms}, refitting after each --- since
\S\ref{sec:mech-built} already shows frozen-parameter ablations understating a
mechanism whose parameters were fitted for a different policy.

\paragraph{L4. Forced-endgame accuracy is still far below the feasible
ceiling.} M2 takes forced-endgame accuracy from \numVfourForcedAcc\% to
\numVfiveForcedAcc\% over \numForcedEightGames{} games per arm, and the measured
ceiling for the best capacity-feasible allocation at those states is
\numFeasibleRight\%. The gap is not closed, and the reason is that feasibility
is exact while the score that ranks the survivors is not: it runs on the
deployed approximate belief. \emph{Resolving experiment:} a calibration study of
M2's allocation estimate against exhaustive enumeration on small reachable
states, to bound the gap between its exact feasibility test and its approximate
score, followed by re-ranking survivors with the exact belief where the state is
small enough to afford it.

\paragraph{L5. A single fitting round.} The configuration evaluated here comes
from \numFitRounds{} round of \numFitGens{} generations
(\S\ref{sec:fitting}); the v0.4 configuration came from
\numVfourFitRounds{} rounds.
The gap between in-panel and held-out performance is already visible in the one
round that was run: the best generation scored \numFitWorst\% against \vfour{}
inside the panel, while the frozen configuration --- the distribution mean the
winner's-curse guard preferred to it (\S\ref{sec:fit-guard}) --- scores
\numPanelVfiveVfour\% against \vfour{} on a bank the fit never saw.
The held-out figure is the one reported, and the difference is a regression to be
expected of a single round rather than an anomaly. \emph{Resolving experiment:}
further rounds on fresh deal banks, each evaluated on a bank the fit never saw,
with the round-by-round in-panel and held-out figures published together so the
regression can be read directly.

\paragraph{L6. The feint exposure is real, replicated, and not addressable by
retuning.} Against the feint archetype \vfive{} scores $\numDecFeintDelta$
points relative to \vfour{}, and does so at both seed banks
(\S\ref{sec:results}). The proximate
cause is identified: the fit raised the weight on ``this player asked here'' from
\numPriorThetaFour{} to \numPriorThetaFive{}, and \S\ref{sec:diag-prior}
establishes that the exposure is over-weighting the policy prior rather than
weighting it at all: against the deceptive archetypes themselves, deleting the
prior outright is worse still, by
\numPriorDeleteCost{} points [\numPriorDeleteCI]. A sweep of that weight over \numThetaSweepN{} values
$\{\numThetaSweep\}$ at two seeds does not identify it: the ordering is unstable
across seeds and every pairwise difference sits inside the cluster-bootstrap
interval, so the fitted value is retained rather than tuned on noise.
\emph{Resolving experiment:} M7, a per-seat online type posterior with
data-biased shrinkage \cite{johanson-dbr}, evaluated against the deception panel
at the seeds used here plus fresh ones, and reported as a frontier rather than an
operating point \cite{johanson-rnash}.

\paragraph{L7. Owner decision D2, partner-aware play, is specified and not
implemented.} The design requires two regimes --- convention-rich,
lower-temperature play with bot teammates, and grounded, less predictable play
with human or unknown teammates --- reported separately, with the self-play
configuration never headlined as the one a human will meet
(\S\ref{sec:mech-partner}). \vfive{} implements neither. It is a single
deterministic policy used in both regimes, and every number in this report is
from the self-play regime. There is no human-teammate evaluation of \vfive{} at
all, which matters most for the one result this study did not obtain by
measurement: the report that a well-trained human outperforms \vfour{} on some
reasoning. \emph{Resolving experiment:} build M6 with the temperature and the
convention flag keyed on the partner regime, then evaluate the two regimes as
separate arms --- bot-teammate self-play, and \vfive{} partnered with a policy it
was not trained with as a stand-in for an unknown teammate --- before any
human-teammate sessions are attempted.

\paragraph{L8. Owner decision D1, conventions behind a flag, has no flag.}
\S\ref{sec:mech-conventions} states the policy: build the signalling machinery
behind an explicit switch, ship it off, and publish the with/without difference
as a result in its own right. The switch does not exist. No codebook, no
convention machinery and no receiver-side decoder were built, so the difference
that was to be published is unmeasured, and the headline configuration is
``conventions off'' only in the trivial sense that there is nothing to turn on.
\emph{Resolving experiment:} implement M5's codebook term behind
\texttt{--conventions}, and report the paired difference between the two settings
on the same deal bank, separately for each partner regime of L7, since a
codebook is worth nothing to a teammate who does not share it.

\subsection{Threats grouped by kind of validity}
\label{sec:lim-validity}

\textbf{Internal validity.} Every ablation in Table~\ref{tab:mechanisms} is a
frozen-parameter ablation: no row was refitted, so no row measures a mechanism's
standalone value under parameters chosen for it. The causal isolation of
\S\ref{sec:diag-features} was performed with a prototype filter applied to the
\vfour{} policy, not with the shipped M1, and the two differ in the fallback rule
and in the ownership scaling. The diagnosis probes are copies of the shipped
policy rather than the shipped policy itself; equality with the original was
checked by reproducing its aggregate statistics to the reported digits after
every instrumentation change, which is a strong check and not a proof of
identity. On the other side of the ledger, the evaluated configuration is the
frozen one: the round-trip assertion that the compiled defaults and the fitted
vector name the same policy passes, and the rules audit reports
\numVerifyViolations{} violations in \numVerifyChecks{} checks with determinism
reproduced.

\textbf{Construct validity.} The provably-dead predicate of
Definition~\ref{def:dead} is a hard deduction and admits no false positives, but
several of the headline statistics that motivate M3 are ground-truth statistics:
the lock census, the ``asks inside a half-suit the team already owns'' rate and
the certificate-ladder measurements all classify half-suits by the true deal, not
by what any observer can establish. \S\ref{sec:diag-information} and
\S\ref{sec:corrections} state the consequence --- the provable version of the
condition holds at \numProvableOwn\% of decisions --- but the reader should not
read those rates as an achievable policy. M2's allocation probability is computed
with the deployed approximate belief, so it is an estimate; the feasibility
constraints it enforces are exact, the score it ranks survivors by is not.
More generally, the posterior throughout this work is exact with respect to the
rules and a uniform prior over consistent deals, which is the policy-agnostic
object; the true posterior weights deals by opponent policy, and the two-scalar
prior examined in \S\ref{sec:diag-prior} is a parametric stand-in for that
correction rather than the correction --- which is also why L6's exposure is a
modelling gap and not a tuning error. Finally, win rate does not capture
declaration reliability, calibration, or the pathology statistics that gate this
work, which is why they are reported separately.

\textbf{External validity.} Every opponent in this study was written inside this
project or ported from an earlier FishBot. The deceptive archetypes of
\S\ref{sec:diag-prior} are commitment deceivers --- they follow a fixed deceptive
rule for the whole game --- and while the withholder is a direct model of the
manoeuvre the project owner reported, it does not reproduce what he actually
did, which was to adapt within a game and across games. The gain of
$\numDecWithholderDelta$ points against it is therefore a result about a fixed
rule, not about a human. The report that a well-trained human outperforms
\vfour{} on some reasoning rests on one player's account of live sessions, is not
a controlled measurement, and is treated here as a hypothesis that generated
experiments rather than as a result. No expert human logs exist for this dialect
and no independently authored engine for it was located, so nothing here tests
the agent against a strategy source outside this project. One point runs the
other way and is worth stating: the three deceptive archetypes, together with
\vtwo{}, the bluffer and the random opponent, were absent from the fitting panel
(\S\ref{sec:fitting}), so the deception panel is a genuinely held-out
measurement of deal bank \emph{and} opponent style, which most of the nine-style
table is not.

\textbf{Statistical validity.} The per-style table is a single seed bank
(\numStyleDeals{} deals in six rotations per cell, seed \numStyleSeed{}), and the
deception panel is two; only the \vfour{} pairing has been run at more
(\numHeadSeeds{} banks, \numHeadDeals{} deals, mean \numHeadWin\%, range
\numHeadRange\%). Several of the diagnosis figures are stable in rate and
unstable in magnitude across seeds, and the pooled values are used above for that
reason. Re-running the information measurement at two fresh seeds reproduced the
rate (\numCertRate\% pooled against \numCertRateOrig\% at the original seed) but
moved the mean improvement from $+\numCertMeanOrig$ to $+\numCertMean$ pooled;
the certificate ladder reached certainty in \numLadderOrigN{} of
\numLadderOrigN{} attempts at the original seed and in \numLadderRate\% pooled,
at a median of \numLadderMedian{} asks rather than two; and the ratio at which a
teammate's best allocation moves toward the truth rather than away from it fell
from $\numMapRatioOrig\!:\!1$ to $\numMapRatioPooled\!:\!1$ pooled. All keep
their sign. The ``information-free at every state'' result of
\S\ref{sec:diag-features} is scoped to \numInfoFreeStates{} states already inside
a repeat cycle and is not a claim about \vfour{}'s asks in general --- indeed
\numOwnLocked\% of all its asks land inside a half-suit its team already owns and
do emit certificates, by accident rather than by choice. Intervals quoted are
deal-clustered percentile bootstrap intervals, which under-cover in evaluations
of this kind \cite{jordan-rleval}.

\textbf{Opponent-population dependence.} Nothing in this report bounds
exploitability, and no such bound is claimed; L2 is the scope statement and
\cite{lisy-lbr} the form the eventual measurement must take. Two of the
mechanisms that would bear on it are unbuilt: while action selection remains a
deterministic argmax, leakage sits at the deterministic bound whatever the
feature weights \cite{alvim-leakage}, so any leak penalty in the ask score is
acting on the wrong variable; and the opponent model remains a fixed global
prior rather than a per-seat one, so the robustness frontier that M7 is supposed
to expose does not yet exist to be reported.

\textbf{Rules-dialect dependence.} All results are for the dialect of
\S\ref{sec:rules}, which is one explicit selection among documented variants
\cite{pagat}. Arbitration by lowest seat is a modelling choice rather than a
rule, and the action cap is a backstop rather than part of the game. The
deadlock measurements in particular are sensitive to the cap only in the patient
configuration, where the cap is what ends the game.

\textbf{Computational and model-class limits.} The ask score remains a
twenty-feature linear function refined by a one-ply expectimax, which caps what
any reweighting can express; M3 and M5 add terms to that class rather than
replacing it. The value model inside it is effectively one feature --- held-out
$R^2$ of \numValueRSq{} for all sixteen against \numValueRSqScore{} for a bias
term plus the score differential alone, where a tree reaches \numValueRSqTree{}
--- so M9's decision to rebuild or retire it is still open, and until it is
taken the declaration rule is scoring on a model whose capacity has been
measured and found wanting. M2's search is exhaustive over at most $3^6$
assignments and returns no candidate at all when an opponent provably holds a
card of the half-suit or when no teammate may hold one, so the declaration path
still needs a defined behaviour at those states. The duplicate-certificate growth
noted in \S\ref{sec:diag-cycle} makes the exact belief progressively more
expensive inside a long cycle without changing what it computes; de-duplication
is an unbuilt engine improvement.

\textbf{Experiments not tied to a single limitation above}, none of which has
been run: a provability oracle over the transportation polytope, to establish how
far the \numProvableOwn\% provable-ownership rate can be widened, which decides
whether M3 can ever gate rather than price; and an adaptive human-like opponent,
or logged human play, to test whether the owner's live result survives a
controlled setting that commitment deceivers do not reproduce.

\textbf{Provenance of the figures in this report.} Of the \numProvUsed{}
distinct quantities quoted in the text, \numProvGenerated{} are emitted directly
from an experiment artifact by \filepath{engine/build_tables_v05.py} and cannot
drift from it; the remaining \numProvTranscribed{} are transcribed out of the
diagnosis reports of \S\ref{sec:diagnosis} rather than regenerated. Transcription
is the weaker guarantee, and it is the mechanism by which a number goes stale
when an artifact is regenerated and the prose is not. It is constrained rather
than eliminated: \filepath{paper/check_provenance.py} requires every transcribed
quantity to sit under a comment naming the artifact it came from and fails the
build if that artifact is absent, and it reports the split above so that this
paragraph cannot itself go stale. Extending the generator to parse the diagnosis
reports would close the gap; it has not been done, and until it is, the
\numProvTranscribed{} transcribed figures carry the same standard of evidence as
the v0.4 study's hand-entered numbers, which is precisely the standard
\S\ref{sec:corrections} criticises elsewhere.
